\documentclass[10pt]{article}

% DOCUMENT LAYOUT
%\usepackage{fullpage}
%\usepackage[cm]{fullpage}
%\usepackage[letterpaper, top=0.5in, bottom=0.5in, left=0.5in, right=0.5in]{geometry} 
%\usepackage[letterpaper]{geometry} 
%\geometry{textwidth=5.7in, textheight=9.0in, marginparsep=7pt, marginparwidth=.6in}
%\geometry{textwidth=6.0in, textheight=9.8in, marginparsep=7pt, marginparwidth=.8in}
%\geometry{textwidth=6in, textheight=9.8in, marginparsep=7pt, marginparwidth=.8in}
\setlength\parindent{0in}

%\addtolength{\oddsidemargin}{-.500in}
%\addtolength{\evensidemargin}{-.500in}
%\addtolength{\textwidth}{2.00in}

%\addtolength{\topmargin}{-.875in}
%\addtolength{\textheight}{1.75in}

% Dynamically control baselinestretch
\usepackage{setspace}

\addtolength{\topmargin}{-1.125in}
\addtolength{\textheight}{2.0in}

\addtolength{\oddsidemargin}{-.375in}
\addtolength{\evensidemargin}{-.375in}
\addtolength{\textwidth}{1.75in}
\addtolength{\marginparsep}{-10pt}
\addtolength{\marginparwidth}{+15pt}

\usepackage{wrapfig}

\usepackage{chngpage}
\usepackage[export]{adjustbox}
\usepackage{mhchem}

\usepackage{sidecap}
\usepackage[abs]{overpic}
\usepackage{wrapfig}

% SPECIAL FORMATTING
% ---- CUSTOM AMPERSAND
\newcommand{\amper}{{\selectfont\itshape\&}}
% ---- MARGIN YEARS
%\newcommand{\years}[1]{\marginpar{\scriptsize #1}}
\newcommand{\years}[1]{\marginpar{\quad \small #1}}
% ---- TALK COUNTER
\newcounter{talknumber}
\setcounter{talknumber}{1}
% ---- ARTICLE COUNTER
\newcounter{articlenumber}
\setcounter{articlenumber}{1}
%\newcommand{\article}{\noindent\marginpar{\scriptsize \arabic{articlenumber}}\addtocounter{articlenumber}{1}}
\newcommand{\article}[2]{
\noindent\marginpar{
  \parbox[c]{0.25in}{\scriptsize \arabic{articlenumber}} 
  \parbox[c]{0.25in}{\scriptsize \href{#2}{#1}}
}\addtocounter{articlenumber}{1}}
% PAPER DESCRIPTION AND THUMBNAIL IMAGE
\newcommand{\desc}[2]{
\noindent\marginpar{
  \parbox[c]{0.50in}{{\includegraphics[width=0.9in]{thumbnails/{#1}}}}
}{\footnotesize\emph{#2}}}
% NEW ARTICLE STYLE WITH IMAGE AND NO NUMBERS
% #1 : thumbnail filename
% #2 : author list
% #3 : article title
% #4 : journal reference
% #5 : DOI, bioRxiv, or ar$\chi$v
% #6 : DOI, bioRxiv, or arXiv link
% #7 : blurb / description
\newcommand{\newarticle}[7]{
\begin{adjustwidth}{-1in}{-1in}  
\begin{tabular}{p{0.9in}p{7in}}
\parbox[c]{0.9in}{\includegraphics[width=0.9in]{thumbnails/#1}} & \parbox[c]{6in}{\setstretch{0.9} {\scriptsize {#2}} \\ {\bf #3}  \\ {\small #4} $\cdot$ \href{#6}{#5} \\ {\footnotesize\emph {#7}}}
\end{tabular}
\end{adjustwidth}
\vspace{0.2in}
}

\newcommand{\newsoftware}[3]{
\begin{adjustwidth}{-1in}{-1in}  
\begin{tabular}{p{0.9in}p{7in}}
\parbox[c]{0.9in}{} & \parbox[c]{6in}{\setstretch{0.9} {\scriptsize {#1}} \\ {\bf #2}  \\ {\footnotesize\emph {#3}}}
\end{tabular}
\end{adjustwidth}
\vspace{0.2in}
}
% ---- TALKS
\newcommand{\talk}[2]{
\noindent\marginpar{
   \scriptsize \scshape
  #2 \\ #1
}}
\newcommand{\numberedtalk}[2]{
\noindent\marginpar{
  \parbox[c]{0.25in}{\scriptsize \arabic{talknumber}} 
% \parbox[c]{0.5in}{\scriptsize {#2}}
  \parbox[c]{0.5in}{\scriptsize \scshape {#1}}
}\addtocounter{talknumber}{1}}

\newcommand{\newtalk}[3]{
\noindent\marginpar{
   \scriptsize \scshape
  #2 \\ #1
}\spaceTalk\\}

% SPACING
\newcommand{\spaceEd}{\vspace{1ex}}
\newcommand{\spaceRes}{\vspace{1ex}}
\newcommand{\spaceFell}{\vspace{0.5ex}}
\newcommand{\spacePub}{\vspace{0.5ex}}
\newcommand{\spaceTalk}{\vspace{1ex}}
\newcommand{\spaceSoc}{\vspace{0.5ex}}
\newcommand{\spacePost}{\vspace{0.5ex}}

% Modify section spacing
\usepackage{titlesec}
\titlespacing*{\section}{0pt}{0.2\baselineskip}{0.2\baselineskip}
\titlespacing*{\subsection}{0pt}{0.2\baselineskip}{0.2\baselineskip}
\titlespacing*{\subsubsection}{0pt}{0.2\baselineskip}{0.2\baselineskip}

% GRAPHICS
\usepackage{graphicx}

% HEADINGS
\usepackage{sectsty} 
\usepackage[normalem]{ulem} 

\usepackage{fancyhdr}
\renewcommand{\headrulewidth}{0pt}
\pagestyle{fancy}
\fancyhf{}
%\fancyfoot[LE,CO]{{\small\textsc{John D. Chodera - Sloan Kettering Institute - \thepage}}}%
\rfoot{\small\textsc{Henrique M. Cezar - Hylleraas Centre - \thepage}}

% Clean up font spacing with micro type
\usepackage{microtype}

% Use a clean, modern-looking font
\usepackage[default]{lato}
% \usepackage{libertine}
\usepackage{FiraSans}
\usepackage[T1]{fontenc}

% Change the appearance of section and subsection headings
\sectionfont{\mdseries\large\underline} 
\subsectionfont{\mdseries\scshape\normalsize} 
\subsubsectionfont{\bfseries\upshape\normalsize} 

% Control tightness of spacing
%\renewcommand{\baselinestretch}{1.00}

% PDF SETUP
% ---- FILL IN HERE THE DOC TITLE AND AUTHOR
\usepackage[bookmarks, colorlinks, breaklinks, pdftitle={Henrique M. Cezar - CV},pdfauthor={Henrique M. Cezar}]{hyperref}  
\hypersetup{linkcolor=blue,citecolor=blue,filecolor=black,urlcolor=blue} 

% Change the typewriter font to something nicer.
\usepackage[ttdefault]{sourcecodepro}

% DOCUMENT
\begin{document}
\reversemarginpar
% FILL IN NAME HERE
{\fontseries{eb}\selectfont \LARGE Henrique Musseli Cezar}\\[0.5cm]

% FILL IN AUTHOR INFORMATION HERE
\begin{minipage}[t]{2.5in}
%\Large CONTACT
%%\includegraphics[width=0.7in,valign=c]{images/john_chodera_sm.pdf}
\hspace{35pt}
\includegraphics[width=1.4in,valign=c]{images/henrique2.jpg}
\end{minipage}
%\quad
\begin{minipage}[t]{3in}
\begin{tabular}{rl}
{\bf email} & \href{mailto:h.m.cezar@kjemi.uio.no}{\href{mailto:h.m.cezar@kjemi.uio.no}{h.m.cezar@kjemi.uio.no}}\\[0.05in]
{\bf github} & \href{https://github.com/hmcezar}{https://github.com/hmcezar}\\[0.05in]
{\bf orcid} & \href{https://orcid.org/0000-0002-7553-0482}{0000-0002-7553-0482}\\[0.05in]
{\bf linkedin} & \href{https://www.linkedin.com/in/hmcezar}{https://www.linkedin.com/in/hmcezar}\\[0.05in]
{\bf mobile} & +47 968 17 973\\[0.05in]
{\bf post} & 
\parbox[t]{3.0in}{
  Hylleraas Centre for Quantum Molecular Sciences\\
  and Department of Chemistry, \\
  University of Oslo, \\
  PO Box 1033 Blindern, 0315 Oslo, Norway}
\end{tabular}
\end{minipage}

\vspace{15pt} 

\noindent\fbox{%
    \parbox{\textwidth}{ 
    {Ph.D in Physics, working with computational physics/chemistry and developing scientific software. 
    My research interests are in the development of new methods for molecular simulation and the use of such methods to solve problems in physical chemistry, biochemistry, and materials science.}
    % My research interests are in the development of new methods for simulation and data analysis, including developing and employing machine learning methods for solving scientific problems.}
    }
}

% \setstretch{1.00} 

\section*{Education and positions}

\noindent\years{2022--}\textbf{Postdoctoral researcher, Hylleraas Centre for Quantum Molecular Sciences, University of Oslo}\\
Further developments in the Hamiltonian hybrid particle-field molecular dynamics and use with metainference
Supervisor: Prof. Michele Cascella \\
\noindent\years{2020--2022}\textbf{Postdoctoral researcher, Institute of Physics, University of S\~ao Paulo}\\
Computational nanofluidics studies on the confinement effects on fluid-solid interfaces.\\
FAPESP Fellowship - Supervisor: Prof. Caetano R. Miranda \\
\noindent\years{2018--2019}\textbf{Postdoctoral researcher, Institute of Physics, University of S\~ao Paulo}\\
Development and applications of Configurational Bias Monte Carlo Method in dyes of interest to organic solar cells \\
CAPES Fellowship - Supervisor: Prof. Kaline R. Coutinho \\
\noindent\years{2015--2018}\textbf{Ph.D. in Physics, Institute of Physics, University of S\~ao Paulo}\\
Implementation and Development of Efficient Algorithm for Intramolecular Deformation with the Monte Carlo Method \\
CNPq Ph.D. Fellowship - Advisor: Prof. Kaline R. Coutinho \\
\noindent\years{2013--2015}\textbf{Master's in Computational Physics, S\~ao Carlos Institute of Physics, University of S\~ao Paulo}\\
Implementation of the Parallel Tempering Monte Carlo Method to the Study of Thermodynamic Properties of Nanoclusters \\
FAPESP Master Fellowship - Advisor: Prof. Juarez L. F. Da Silva \\
\noindent\years{2009--2012}\textbf{B.S. in Computational Physics, S\~ao Carlos Institute of Physics, University of S\~ao Paulo}\\
Undergraduate research entitled: Fluid Dynamics in Porous Media with Lattice Boltzmann \\ 
CNPq Scientific Initiation Fellowship - Advisor: Prof. Leonardo P. Maia

% \section*{Complementary Education}
% \noindent\years{2015}{\textbf{Coursera Online Course:} Algorithms: Design and Analysis, Part 1. (Duration: 36h). Stanford University.} \\
% \noindent\years{2006--2008}{\textbf{Technical Course:} Informatics. Col\'egio Divino Salvador.}


\section*{Approved grants}
\noindent\years{2020}{\textbf{FAPESP} Two-years Postdoc Fellowship.} \\
\noindent\years{2013}{\textbf{FAPESP} Two-years Masters Fellowship.}

% \section*{Fellowships}

% \noindent\years{2015--}CNPq Ph.D. Fellowship\\
% \noindent\years{2013--2015}FAPESP Master Fellowship\\
% \noindent\years{2012--2012}CNPq Scientific Initiation Fellowship

\section*{Awards}
\noindent\years{2022}{\textbf{Best PhD thesis} in the field of atomic and molecular physics, by the Brazilian Physical Society.} \\
\noindent\years{2022}{\textbf{Honorable mention} at the annual prize of the Institue of Physics, University of S\~ao Paulo, at the ``Highlight in instrumentation and software development'' category.} \\
\noindent\years{2020}{\textbf{Top Downloaded Paper}: ``Solvent effect on the syn/anti conformational stability: A comparison between conformational bias Monte Carlo and molecular dynamics methods'', Wiley.} \\
\noindent\years{2019}{\textbf{3rd place} in the best flash talks award at the XX Simp\'osio Brasileiro de Qu\'imica Te\'orica.} \\
\noindent\years{2017}{\textbf{Best poster} award at the XIX Simp\'osio Brasileiro de Qu\'imica Te\'orica.} \\
\noindent\years{2017}{\textbf{Among the 12 finalists} of the ``My Thesis in 180 Seconds'', organized by Swissnex Brazil.} \\
\noindent\years{2015}{\textbf{Honorable mention} at the best poster award of the II Workshop on Biomolecular Theory-Experiment Interplay.}

\section*{Teaching Experience}
\noindent\years{2019}{\textbf{F\'isica II (Basic Physics II)}: Teacher for the second semester physics class for the mechanical engineering majors at the Escola Polit\'ecnica -- University of S\~ao Paulo.} \\
\noindent\years{2019}{\textbf{F\'isica I (Basic Physics I)}: Teacher for the first semester physics class for the mechanical engineering majors at the Escola Polit\'ecnica -- University of S\~ao Paulo.} \\
\noindent\years{2017}{\textbf{F\'isica Moderna I (Modern Physics I)}: Teaching assistant for the introductory modern physics class for the physics majors of the Institute of Physics -- University of S\~ao Paulo, under the supervision of Prof. Dr. Sylvio Canuto.}

\section*{Research Experience}
My research is mainly focused on the development of algorithms for classical molecular simulation. 
Specifically, I have experience implementing and using Monte Carlo methods and molecular dynamics for the simulation of solute-solvent, interfacial and biomolecular systems.
At the Hylleraas Centre, I have been maintaining Hylleraas MD, a Hamiltonian hybrid particle-field software, in which I have implemented a interface with PLUMED.
I have also given contributions to PLUMED development regarding the implementation of small angle neutron scattering (SANS) and the Python MPI interface.
On top of the molecular simulation methods, I have some experience with global optimization algorithms, clustering algorithms, fitting of interatomic potentials, and density functional theory.

Recently, I have been employing metainference, a Bayesian inference method, to bias coarse-grain simulations towards agreement with small angle scattering spectra in collaboration with experimentalists.
With my collaborators, we are exploring the use of machine learning potentials (NequIP and Allegro) and their integration with the the Hylleraas Software Platform.
I'm also leading a project to develop a method for the classification of structures in solute-solvent systems (ClustTraj).

\section*{Languages}
Portuguese (native), English (fluent) and Norwegian (novice).

\section*{Technical Skills}
I am proficient with Python, Fortran, C, Bash, and C++, and have experience with OpenMP and MPI parallelization, code profiling, and optimization. 
My Python projects have used libraries like NumPy, Pandas, SciPy, mpi4py, scikit-learn and matplotlib.
I also have some experience using GitHub actions for automating unit testing, deploying documentation, and publishing Python packages.
I know how to use HPC systems like SLURM and PBS, and have experience with Docker and Singularity containers.

Concerning scientific software, I am familiar with DICE, GOTNano, HylleraasMD, Cassandra, GROMACS, and LAMMPS for molecular simulations. 
For electronic structure calculations, I have some experience with Gaussian, FHI-AIMS and Quantum Espresso.

% \setstretch{1.00} 

\eject

%%%%%%%%%%%%


\section*{Software development}
\href{https://github.com/hmcezar}{GitHub}, 
\href{https://gitlab.com/hmcezar}{GitLab}

Apart from many scripts and smaller software (see the links above), my major contributions are developing:\\

\newsoftware{\textbf{Cezar HM}, Canuto S, Coutinho K}{DICE}{\href{https://portal.if.usp.br/dice}{https://portal.if.usp.br/dice} \\ 
A Monte Carlo code for the molecular simulation of liquids and gases using classical force fields. My main contribution is the coding and development of a Configurational Bias Monte Carlo method used to sample the internal degrees of freedom of the molecules, and parallelization. The code is written in Fortran and has OpenMP parallelization.}

\newsoftware{Rondina GG, \textbf{Cezar HM}, Da Silva JLF}{GOTNano}{ 
A code for global optimization and sampling of thermodynamic properties of nanoclusters. My main contribution is the implementation of the Parallel Tempering Monte Carlo algorithm used to sample the thermodynamic properties. The code is written in C and has OpenMP parallelization.}

\newsoftware{Ledum M, Carrer M, Sen S, Li X, \textbf{Cezar HM}, Cascella M, Bore SL}{Hylleraas MD}{\href{https://github.com/Cascella-Group-UiO/HyMD}{https://github.com/Cascella-Group-UiO/HyMD} \\ 
A molecular dynamics software that implements the Hamiltonian hybrid particle-field method. I am currently the main maintainer of the software, and my main contribution is the PLUMED interface. HyMD is written in Python and Fortran and uses MPI for parallelization.}

\newsoftware{\textbf{Cezar HM}, Ribeiro RB}{ClustTraj}{\href{https://github.com/hmcezar/clusttraj}{https://github.com/hmcezar/clusttraj} \\ 
A parallel Python package that performs clustering over an atomistic molecular simulation trajectory, searching for the minimum RMSD between each structure and classifying using hierarchical clustering, with options to deal with solute solvent systems.}

\newsoftware{\textbf{Cezar HM}, Sveinsson HA, Bore SL \textit{et al.}}{Hylleraas Software Platform NequIP/Allegro interface}{\href{https://gitlab.com/hylleraasplatform/hylleraas-interfaces}{https://gitlab.com/hylleraasplatform/hylleraas-interfaces} \\ 
Interface in the Hylleraas Software Platform for the training and evaluation of NequIP/Allegro models. The interface is used in the active learning package under development in the same platform. I highlight that the interface for other software in the Git repository are contributions from other people.}

% \newsoftware{\textbf{Cezar HM}}{DICEtools}{\href{https://github.com/hmcezar/dicetools}{https://github.com/hmcezar/dicetools} \\ 
% A package containing several scripts used to prepare DICE inputs and analyze simulation data.}

% \newsoftware{\textbf{Cezar HM}}{pdb2lmp}{\href{https://github.com/Sampa-USP/pdb2lmp}{https://github.com/Sampa-USP/pdb2lmp} \\ 
% Create LAMMPS topologies from structure files.}

% \newsoftware{\textbf{Cezar HM}}{nanotube-tools}{\href{https://github.com/Sampa-USP/nanotube-tools}{https://github.com/Sampa-USP/nanotube-tools} \\ 
% Set of tools to manipulate, create and characterize nanotubes.}

%%%%%%%%%%%%%%%%%%%%%%%%%%%%%%%%%%%%%%%%%%%%%%%%%%%%%%%%%%%%%%%%%%%%%%%%%%%%%%%%
% PUBLICATIONS
%%%%%%%%%%%%%%%%%%%%%%%%%%%%%%%%%%%%%%%%%%%%%%%%%%%%%%%%%%%%%%%%%%%%%%%%%%%%%%%%

\section*{Publications}

%%%%%%%%%%%%

\textit{h-index: 8, citations: 171 (August 9th, 2024)} \\
\href{https://www.webofscience.com/wos/author/record/445303}{Web of Science}, 
\href{https://www.scopus.com/authid/detail.uri?authorId=57188802591}{Scopus},
\href{https://orcid.org/0000-0002-7553-0482}{ORCID}, 
\href{https://scholar.google.com.br/citations?hl=pt-BR&user=LtBk3gEAAAAJ}{Google Scholar} 

% \subsection*{In Press}

%\eject

%%%%%%%%%%%%

\subsection*{Published}

\newarticle{diff_hymd.pdf}{Carrer M, \textbf{Cezar HM}, Bore SL, Ledum M, Cascella M}{Learning Force Field Parameters from Differentiable Particle-Field Molecular Dynamics}{\emph{Journal of Chemical Information and Modeling} 64:5510, 2024}{DOI}{https://doi.org/10.1021/acs.jcim.4c00564}{We develop a differentiable molecular dynamics software based on the Hamiltonian hybrid particle-field formalism, and use it to establish a protocol for automated optimization of force field parameters.}

\newarticle{cnt_gas_potential.png}{\textbf{Cezar HM}, Lanna TD, Damasceno DA, Kirch A, Miranda CR}{Revisiting greenhouse gases adsorption in carbon nanostructures: advances through a combined first-principles and molecular simulation approach}{\emph{Applied Surface Science} 671:160659, 2024}{DOI}{https://doi.org/10.1016/j.apsusc.2024.160659}{We investigate the adsorption properties of \ce{CH4} and \ce{CO2} in carbon nanotubes (CNTs) using different force fields, showing that the adsorption energies are underestimated. A new interatomic potential is proposed based on the first-principles calculations.}

\newarticle{non-bo-structure.png}{Lang L, \textbf{Cezar HM}, Adamowicz L, Pedersen TB}{Quantum Definition of Molecular Structure}{\emph{Journal of the American Chemical Society} 146:1760, 2024}{DOI}{https://doi.org/10.1021/jacs.3c11467}{In this manuscript we use Monte Carlo to sample an all-particle wavefunction, and propose a new definition of molecular structure beyond the Born-Oppenheimer approximation.}

\newarticle{sans-plumed.jpeg}{\textbf{Cezar HM}, Cascella M}{SANS Spectra with PLUMED: Implementation and Application to Metainference}{\emph{Journal of Chemical Information and Modeling} 63:4979, 2023}{DOI}{https://dx.doi.org/10.1021/acs.jcim.3c00724}{In this Application Note we describe our implementation of Small Angle Neutron Scattering spectra including the use of a resolution function and its application to metainference simulations.}

\newarticle{flip-flop.jpeg}{Carrer M, Nielsen JE, \textbf{Cezar HM}, Lund R, Cascella M, Soares TA}{Accelerating Lipid Flip-Flop at Low Concentrations: A General Mechanism for Membrane Binding Peptides}{\emph{The Journal of Physical Chemistry Letters} 14:7014, 2023}{DOI}{https://doi.org/10.1021/acs.jpclett.3c01284}{We show that the flip-flop of lipids in a membrane can be accelerated by the presence of peptides. The mechanism is entropic, due to the restriction in mobility of lipids introduced by the presence of the peptide, and is not necessarily driven by smaller enthalpic barriers.}

\newarticle{dpbimH.png}{Franco LR, Toledo KCF, Matias TA, Benavides PA, \textbf{Cezar HM}, Araujo CMG, Coutinho K, Araki K}{Unraveling the acid-base characterization and solvent effects on structural and electronic properties of a bis-bidentated bridging ligand}{\emph{Physical Chemistry Chemical Physics} 24:10222, 2022}{DOI}{https://doi.org/10.1039/D1CP03912A}{We report a combined theoretical-experimental study of the bis(2-pyridyl)-bis-imidazole in solution, determining the protonation sites, most favorable conformations and absorption spectra in different solvents.}

\newarticle{dha-vhf.png}{Cardenuto MH, \textbf{Cezar HM}, Mikkelsen KV, Sauer SPA, Coutinho K, Canuto S}{A QM/MM study of the conformation stability and electronic structure of the photochromic switches derivatives of DHA/VHF in acetonitrile solution}{\emph{Spectrochimica Acta Part A: Molecular and Biomolecular Spectroscopy} 251:119434, 2021}{DOI}{https://doi.org/10.1016/j.saa.2021.119434}{The absorption spectra of photoswichtes composed of two units of photochromic molecules, namely, dihydroazulene (DHA)/vinylheptafulvene(VHF) is studied including explicit solvation and considering the rotamers that are exhibited at room temperature. We obtain the free energies and perform clustering analysis, rationalizing the absorption spectra at the different solvents.}

\newarticle{cbmc_dice.png}{\textbf{Cezar HM}, Canuto S, Coutinho K}{DICE: A Monte Carlo code for molecular simulation including Configuration Bias Monte Carlo method}{\emph{Journal of Chemical Information and Modeling} 60:3472, 2020}{DOI}{https://doi.org/10.1021/acs.jcim.0c00077}{We describe DICE, a software that uses Monte Carlo methods to perform molecular simulation, focusing on solute-solvent systems. We introduce our implementation of the Configurational Bias Monte Carlo method that includes new methodological developments and employ the method to study different molecules such as alkanes, 1,2-dichloroethane and a large boron phtalocyanine.}

\newarticle{mox_spectroscopy.png}{\textbf{Cezar HM}, Canuto S, Coutinho K}{Understanding the Absorption Spectrum of Mesityl Oxide Dye in Solvents of Different Polarities}{\emph{Journal of Molecular Liquids} 307:112924, 2020}{DOI}{https://doi.org/10.1016/j.molliq.2020.112924}{We investigate the conformational and solvent effects on the UV/vis spectroscopy of the mesityl oxide dye. We show that a correct description of the \textit{syn}/\textit{anti} population and the use of explicit solvent molecules is fundamental to describe the solvatochromic shifts.}

\newarticle{ptmc_pt_alloys.png}{\textbf{Cezar HM}, Rondina GG, Da Silva JLF}{Thermodynamic Properties of 55-Atom Pt-based Nanoalloys: Phase Changes and Structural Effects on the Electronic Properties}{\emph{Journal of Chemical Physics} 151:204301, 2019}{DOI}{http://dx.doi.org/10.1063/1.5125689}{We investigate the effects of temperature on the structure of Pt-based nanoalloys, performing Parallel Tempering Monte Carlo simulations using many-body potentials and DFT calculations to see how the structural changes affect the density of states.}

\newarticle{cbmc_mox.pdf}{\textbf{Cezar HM}, Canuto S, Coutinho K}{Solvent effect on the syn/anti conformational stability: A comparison between conformational bias Monte Carlo and molecular dynamics methods}{\emph{International Journal of Quantum Chemistry} 119:e25688, 2019}{DOI}{http://dx.doi.org/10.1002/qua.25688}{We compare the Configurational Bias Monte Carlo (CBMC) and molecular dynamics (MD) sampling of the mesityloxide molecule in gas phase and solution, showing that that while CBMC performs an ergodic sampling, MD has troubles overcoming the 10 kcal/mol energy barrier between the \textit{syn} and \textit{anti} conformations.}

\newarticle{similarity_lj.pdf}{\textbf{Cezar HM}, Rondina GG, Da Silva JLF}{Parallel Tempering Monte Carlo Combined with Clustering Euclidean Metrics Analysis to Study the Thermodynamic Stability of Lennard-Jones Nanoclusters}{\emph{Journal of Chemical Physics} 146:064114, 2017}{DOI}{http://dx.doi.org/10.1063/1.4975601}{We use a clustering algorithm to analyze trajectories from Parallel Tempering Monte Carlo simulations of Lennard-Jones nanoclusters, identifying the phase changes and most most frequent structures with great accuracy using a very simple approach.}

\newarticle{alloys_douglas.pdf}{De Souza DG, \textbf{Cezar HM}, Rondina GG, De Oliveira MF, Da Silva JLF}{A Basin-hopping Monte Carlo Investigation of the Structural and Energetic Properties of 55- and 561-atom Bimetallic Nanoclusters: the Examples of the ZrCu, ZrAl, and CuAl Systems}{\emph{Journal of Physics: Condensed Matter} 28:175302, 2016}{DOI}{http://dx.doi.org/10.1088/0953-8984/28/17/175302}{We report a basin-hopping Monte Carlo investigation of the structural and energetic properties of bimetallic ZrCu, ZrAl, and CuAl nanoclusters with 55 and 561 atoms, showing how the atoms of the different species are distributed in the structure, and a trend of more spherical structures at 50-50 compositions.}

% \subsection*{Accepted for publication}

\subsection*{Preprints}
\newarticle{sio2_nts.png}{\textbf{Cezar HM}, Miranda CR}{Water adsorption in ultrathin silica nanotubes}{arXiv}{DOI}{https://doi.org/10.48550/arXiv.2307.09965}{We investigate the structure and adsorption properties for three sizes of \ce{SiO2} nanotubes, understanding the adsorption in terms of the accessible volume.}

\newarticle{n2-isotherms.png}{Martins Junior CA, \textbf{Cezar HM}, Damasceno DA, Miranda CR}{Effects of van der Waals interaction on the \ce{N2} adsorption on carbon nanotubes: proposal of new force field parameters}{arXiv}{DOI}{https://doi.org/10.48550/arXiv.2307.10117}{We investigate the adsorption of \ce{N2} in CNTs, proposing an interatomic potential that correctly describe the energies for different adsorption sites, leading to a more accurate description of adsorption isotherms.}

\newarticle{cnt-mec-prop.png}{Damasceno DA, \textbf{Cezar HM}, Lanna TD, Kirch A, Miranda CR}{Mechanical and adsorption properties of greenhouse gases filled carbon nanotubes}{arXiv}{DOI}{https://doi.org/10.48550/arXiv.2307.11708}{Using molecular simulations we investigate the mechanical stability of CNTs filled with mixtures of \ce{CO2} and \ce{CH4}, finding that under compression several parameters can play a role in the in the materials response.}

% this is not in a preprint repository yet
% \newarticle{nbzr_stbs.png}{Salvador CAF, \textbf{Cezar HM}, Miranda CR}{Atomistic simulations of symmetric tilt grain boundaries in Nb-1Zr alloys}{}{}{}{We investigate the mechanical properties of three different symmetric tilt grain boundaries (STBs) at the \ce{Nb-1Zr} at.\% compositions. Using non equilibrium molecular dynamics simulations and Monte Carlo moves, we obtain that the \ce{Zr} stays at the grain boundray, and the addition of small amounts of \ce{Zr} can help pin the STBs. We also analyze the stress-strain curves and fracture patterns.}

\section*{Selected invited and contributed talks}
\talk{Jul 2024}{Dublin, Ireland}\textbf{Metainference simulations to interpret small angle scattering experiments} \\ 9th EuChemS Chemistry Congress\\[0.05in]

\talk{Sep 2019}{S\~ao Carlos, Brazil}\textbf{DICEtools and Clustering-Traj: From simulation setup to data analysis} \\ III New Energies Innovation Center - Computational Material Science Division Workshop\\[0.05in]

\talk{Feb 2019}{S\~ao Carlos, Brazil}\textbf{Monte Carlo Methods} \\ II New Energies Innovation Center - Computational Material Science Division Workshop (Invited talk)\\[0.05in]

\talk{Sep 2018}{Bras\'ilia, Brazil}\textbf{Comparison Between Configurational Bias Monte Carlo and Molecular Dynamics Sampling for Sampling Conformational Stability with High Energy Barriers} \\ VII SeedMol  \\[0.05in]

\talk{Aug 2017}{Livorno, Italy}\textbf{Configurational Bias Monte Carlo of Molecules in Solvent and Comparison with Molecular Dynamics} \\ Coding Solvation Workshop \\[0.05in]

% \talk{May 2017}{Recife, PE}\textbf{Comparsion of the sampling of Molecular Dynamics and Configurational Bias Monte Carlo for 1,2- dichloroethane and octane.} \\ III Advanced School on Biomolecular Simulation \\[0.05in]

\talk{Mar 2017}{Rio de Janeiro, Brazil}\textbf{How Indiana Jones is actually a bad explorer when it comes to molecules and how I can do better} \\ Finals of My Thesis in 180 seconds (Swissnex
Brazil) \\[0.05in]

\talk{Sep 2014}{Jo\~ao Pessoa, Brazil}\textbf{Parallel Tempering Monte Carlo Investigation of Phase-Changes in Nanoclusters} \\ XIII Brazilian Materials Research Society Meeting \\[0.05in]


% \section*{Poster presentation}
% \talk{Feb 2021}{Online}\textbf{Revisiting greenhouse gases adsorption in carbon nanotubes:
% advances by ab initio optimized potentials} \\ 20th International Workshop on Computational Physics and Materials Science: Total Energy and Force Methods \\[0.05in]

% \talk{Nov 2019}{Jo\~ao Pessoa, PB}\textbf{Conformational effects on the solvatochromism of mesityl oxide: the importance of an ergodic sampling} \\ XX SBQT \\[0.05in]

% \talk{Feb 2019}{S\~ao Carlos, SP}\textbf{Application of Configurational Bias Monte Carlo method for the simulation of dyes and porous nanoparticles} \\ II New Energies Innovation Center - Computational Material Science Division Workshop \\[0.05in]

% \talk{Oct 2018}{S\~ao Carlos, SP}\textbf{Implementation of Configurational Bias Monte Carlo method for the simulation of flexible molecules in solution} \\ I New Energies Innovation Center - Computational Material Science Division Workshop \\[0.05in]

% \talk{Nov 2017}{\'Aguas de Lind\'oia, SP}\textbf{Implementation of Configurational Bias Monte Carlo Method to Sample Flexible Solute in Solvent Media} \\ XIX SBQT \\[0.05in]

% \talk{Aug 2017}{Munich, Germany}\textbf{Configurational Bias Monte Carlo method to sample molecular flexibility: The case of octane and 1,2-dichloroethane} \\ WATOC 2017 \\[0.05in]

% \talk{Aug 2015}{Maresias, SP}\textbf{Parallel Tempering Monte Carlo applied to the study of transition-metal nanoclusters} \\ II Workshop on Biomolecular Theory-Experiment Interplay \\[0.05in]

% \talk{Aug 2013}{Trieste, Italy}\textbf{Thermodynamic properties of nanoclusters: An investigation with Parallel Tempering Monte Carlo} \\ Density Functional Theory and Beyond: Computational Materials Science for Real Materials \\[0.05in]

% \talk{Oct 2012}{S\~ao Paulo, SP}\textbf{Estudo de din\^amica de fluidos em meios porosos com m\'etodos de Boltzmann na rede} \\ XX SIICUSP \\

\section*{Conferences organized}
\talk{May 2017}{Recife, Brazil}\textbf{III Advanced School on Biomolecular Simulation} \\ Member of the local committee.

% \eject

\vfill{}
% \hrulefill
% credits to Chodera's template
% \begin{center}
% {\scriptsize
% The \LaTeX{} template for this CV is forked from \href{https://github.com/jchodera/latex-cv}{https://github.com/jchodera/latex-cv}}
% \end{center}

\end{document}
